%% kapa: exam=1 round=36 year=2025 period=1 form=A subject=real_estate subject_ko=부동산학원론 seq=24 number=104
%% source: Q-Net 감정평가사 기출문제 (cst00309) artlSeq=5231525 file=raw/1008/5231525/2025년 제36회 감정평가사 1차시험 1교시 문제지 원본.pdf
%% conversion: kordoc 4.24.51 -> tools/gen1.py (auto); review=proofread-sonnet; flags=none
\begin{question}{104}
도시 및 부동산개발에 관한 설명으로 옳지 않은 것은?
\begin{choices}
  \item ｢도시 및 주거환경정비법｣상 “재개발사업”이란 정비기반시설이 열악하고 노후ㆍ불량 건축물이 밀집한 지역에서 주거환경을 개선하거나 상업지역ㆍ공업지역 등에서 도시기능의 회복 및 상권활성화 등을 위하여 도시환경을 개선하기 위한 사업을 말한다.
  \item ｢개발이익 환수에 관한 법률｣상 “개발이익”이란 개발사업의 시행이나 토지이용계획의 변경, 그 밖에 사회적ㆍ경제적 요인에 따라 정상지가(正常地價)상승분을 초과하여 개발사업을 시행하는 자나 토지 점유자에게 귀속되는 토지 가액의 증가분을 말한다.
  \item ｢국토의 계획 및 이용에 관한 법률｣상 “공동구”란 전기ㆍ가스ㆍ수도 등의 공급설비, 통신시설, 하수도시설 등 지하매설물을 공동 수용함으로써 미관의 개선, 도로구조의 보전 및 교통의 원활한 소통을 위하여 지하에 설치하는 시설물을 말한다.
  \item ｢부동산개발업의 관리 및 육성에 관한 법률｣상 “부동산개발업”이란 타인에게 공급할 목적으로 부동산개발을 수행하는 업을 말한다.
  \item ｢도시개발법｣상 “도시개발사업”이란 도시개발구역에서 주거, 상업, 산업, 유통, 정보통신, 생태, 문화, 보건 및 복지 등의 기능이 있는 단지 또는 시가지를 조성하기 위하여 시행하는 사업을 말한다.
\end{choices}
\end{question}
\answer{2}
